\subsection{Álgebras finitas y étales.}
\label{subsection:IV.18.3}

\begin{proposition}[18.3.1]
\label{IV.18.3.1}
Sean $A$ un anillo y $B$ una $A$-álgebra de presentación finita.
\begin{enumerate}
  \item[(i)] Para que $B$ sea una $A$-álgebra no ramificada, es necesario y suficiente que $B$ sea un $(B\otimes_A B)$-módulo proyectivo.
  \item[(ii)] Supongamos además que $B$ es una $A$-álgebra finita. Para que $B$ sea una $A$-álgebra étale, es necesario y suficiente que $B$ sea un $A$-módulo proyectivo y un $(B\otimes_A B)$-módulo proyectivo.
\end{enumerate}
\end{proposition}

Por supuesto, la estructura de $(B\otimes_A B)$-módulo sobre $B$ es la que proviene de la estructura de $(B\otimes_A B)$-álgebra sobre $B$ correspondiente al $A$-homomorfismo canónico de anillos $B\otimes_A B\to B$, que es sobreyectivo y tiene núcleo $\mathfrak J_{B/A}$ (0, 20.4.1).

(i) Decir que el morfismo $\Spec(B)\to\Spec(A)$ es localmente de presentación finita equivale a decir que $B$ es una $A$-álgebra de presentación finita (1.4.6). Decir

\oldpage[IV-4]{115}

que $B$ es una $A$-álgebra no ramificada significa entonces (17.4.2) que $\Spec((B\otimes_A B)/\mathfrak J_{B/A})$ es un subesquema inducido sobre un subconjunto \emph{abierto y cerrado} de $\Spec(B\otimes_A B)$, y se sabe que, para que esto ocurra, es necesario y suficiente que $\mathfrak J_{B/A}$ sea un ideal \emph{factor directo} de $B\otimes_A B$ (Bourbaki, \emph{Alg. comm.}, cap. II, \textsection~4, n.\textsuperscript{o}~3, prop. 15); pero equivale a decir que el $(B\otimes_A B)$-módulo cociente $(B\otimes_A B)/\mathfrak J_{B/A}$ es proyectivo (Bourbaki, \emph{Alg.}, cap. II, 3.\textsuperscript{a} ed., \textsection~2, n.\textsuperscript{o}~2, prop. 4).

(ii) Si se recuerda que un $A$-módulo plano de presentación finita es proyectivo y recíprocamente (Bourbaki, \emph{Alg. comm.}, cap. II, \textsection~5, n.\textsuperscript{o}~2, cor. 2 del th. 1), la afirmación de (ii) se sigue de la de (i) y de (17.6.2).

\begin{proposition}[18.3.2]
\label{IV.18.3.2}
Sean $A$ un anillo y $\mathfrak J$ un ideal de $A$ tal que, para la topología $\mathfrak J$-preádica, $A$ sea separado y completo; pongamos $A_0=A/\mathfrak J$. Entonces el funtor
\[
  B\longmapsto B\otimes_A A_0
\]
es una equivalencia de la categoría de las $A$-álgebras finitas y étales con la categoría de las $A_0$-álgebras finitas y étales.
\end{proposition}

Demostraremos primero el lema siguiente:

\begin{lemma}[18.3.2.1]
\label{IV.18.3.2.1}
Sean $A$ un anillo y $\mathfrak J$ un ideal de $A$ tal que, para la topología $\mathfrak J$-preádica, $A$ sea separado y completo.
\begin{enumerate}
  \item[(i)] Todo $A$-módulo proyectivo de tipo finito $M$ es separado y completo para la topología $\mathfrak J$-preádica y, por tanto, límite proyectivo de los $(A/\mathfrak J^{n+1})$-módulos proyectivos $M/\mathfrak J^{n+1}M$.
  \item[(ii)] Recíprocamente, pongamos $A_n=A/\mathfrak J^{n+1}$, y sea $(M_n)$ un sistema proyectivo de $A_n$-módulos tal que, para todo $n$, el homomorfismo $M_{n+1}\otimes_{A_{n+1}}A_n\to M_n$ deducido del homomorfismo de transición $M_{n+1}\to M_n$ sea biyectivo. Supongamos además que los $M_n$ son proyectivos y que $M_0$ es de tipo finito. Entonces $M=\varprojlim M_n$ es un $A$-módulo proyectivo de tipo finito tal que el homomorfismo canónico $M\otimes_A A_0\to M_0$ es biyectivo.
\end{enumerate}
\end{lemma}

(i) Existe un $A$-módulo libre de tipo finito $L$ tal que $M$ es isomorfo a un sumando directo de $L$; como $L$ es separado para la topología $\mathfrak J$-preádica, lo mismo ocurre con todo submódulo $N$ de $L$, puesto que $\mathfrak J^{n+1}N\subset\mathfrak J^{n+1}L$; en particular, $M$ es separado para esta topología y, como el homomorfismo sobreyectivo $f:L\to M$ es continuo para la topología $\mathfrak J$-preádica, su núcleo $N$ es cerrado para la topología inducida por la de $L$; puesto que $L$ es completo y $f$ un morfismo estricto, se concluye que $M$ es completo (Bourbaki, \emph{Top. gén.}, cap. IX, 2.\textsuperscript{a} ed., \textsection~3, n.\textsuperscript{o}~1, prop. 4).

(ii) Se sigue del lema de Nakayama que, si $M_0$ está generado por una familia finita $(x_{i0})$ de $r$ elementos y si, para todo $n\geq 1$, $x_{in}$ es un elemento de $M_n$ del que $x_{i,n-1}$ es la imagen en $M_{n-1}$, entonces $(x_{in})$ $(1\leq i\leq r)$ es un sistema de generadores de $M_n$ (Bourbaki, \emph{Alg. comm.}, cap. II, \textsection~3, cor. 2 de la prop. 4). Dicho esto, para todo $n$, pongamos $L_n=A_n^r$; si $(e_{in})_{1\leq i\leq r}$ es la base canónica de $L_n$, sea $u_n:L_n\to M_n$ la aplicación $A_n$-lineal tal que $u_n(e_{in})=x_{in}$ para todo $i$. Por hipótesis, se tiene una sucesión exacta escindida
\[
  0\longrightarrow N_n\xrightarrow{v_n}L_n\xrightarrow{u_n}M_n\longrightarrow 0
\]

\oldpage[IV-4]{116}

y, como $L_n=L_{n+1}/\mathfrak J^{n+1}L_{n+1}$ y $M_n=M_{n+1}/\mathfrak J^{n+1}M_{n+1}$, las flechas verticales del diagrama conmutativo
\[
\xymatrix{
0\ar[r] & N_{n+1}\ar[r]^{v_{n+1}}\ar[d] & L_{n+1}\ar[r]^{u_{n+1}}\ar[d] & M_{n+1}\ar[r]\ar[d] & 0\\
0\ar[r] & N_n\ar[r]_{v_n} & L_n\ar[r]_{u_n} & M_n\ar[r] & 0
}
\]
son las tres sobreyectivas. Ahora bien, se tiene $M=\varprojlim M_n$ y $L=A^r=\varprojlim L_n$; si se pone $N=\varprojlim N_n$, $u=\varprojlim u_n$, $v=\varprojlim v_n$, se tiene, en virtud de ($0_{\mathrm{III}}$, 13.2.2), la sucesión exacta
\begin{equation}
  0\longrightarrow N\xrightarrow{v}L\xrightarrow{u}M\longrightarrow 0.
  \tag{18.3.2.2}\label{IV.18.3.2.2}
\end{equation}
Además, como para todo $n$, $v_n$ es invertible por la izquierda y $M_n$ es un $A_n$-módulo proyectivo, se sigue de (0, 19.1.8) que la sucesión exacta (18.3.2.2) es escindida, lo que demuestra el lema.

Dicho esto, demostremos primero que el funtor del enunciado de (18.3.2) es \emph{plenamente fiel}. Pongamos, como en el lema, $A_n=A/\mathfrak J^{n+1}$; sean $B$, $C$ dos $A$-álgebras finitas y étales, y pongamos, para todo $n$, $B_n=B\otimes_A A_n$, $C_n=C\otimes_A A_n$; en virtud de (18.3.1) y (18.3.2.1), $B$ y $C$ son separadas y completas para la topología $\mathfrak J$-preádica, y se tiene $B=\varprojlim B_n$, $C=\varprojlim C_n$; además, todo homomorfismo de $A$-álgebras $u:B\to C$ es continuo para las topologías $\mathfrak J$-preádicas y proporciona, por tanto, un sistema proyectivo de homomorfismos de $A_n$-álgebras $u_n=u\otimes 1:B_n\to C_n$, del que es el límite proyectivo; siendo evidente la recíproca, se tiene así una biyección canónica
\[
  \Hom_{A\text{-alg.}}(B,C)\xrightarrow{\sim}\varprojlim\Hom_{A_n\text{-alg.}}(B_n,C_n).
\]
Pero como $B$ y $C$ son $A$-álgebras étales, se sigue también de inmediato de (18.1.2) que la aplicación canónica
\[
  \Hom_{A_{n+1}\text{-alg.}}(B_{n+1},C_{n+1})\longrightarrow\Hom_{A_n\text{-alg.}}(B_n,C_n)
\]
es biyectiva para $n\geq 0$, lo que termina de demostrar que la aplicación canónica
\[
  \Hom_{A\text{-alg.}}(B,C)\longrightarrow\Hom_{A_0\text{-alg.}}(B_0,C_0)
\]
es biyectiva.

Para terminar de demostrar (18.3.2), basta ver que para toda $A_0$-álgebra finita y étale $B_0$, existe una $A$-álgebra finita y étale $B$ y un $A_0$-isomorfismo $B_0\xrightarrow{\sim}B\otimes_A A_0$. Ahora bien, se sigue de (18.1.2) que existe un sistema proyectivo $(B_n)$ tal que $B_n$ es una $A_n$-álgebra étale y los homomorfismos $B_{n+1}\otimes_{A_{n+1}}A_n\to B_n$ son biyectivos. Se sigue de (18.3.1) y (18.3.2.1) que la $A$-álgebra $B=\varprojlim B_n$ es un $A$-módulo proyectivo de tipo finito y que $B_0$ es isomorfa a $B\otimes_A A_0$. Para demostrar que $B$ es una $A$-álgebra étale, basta, en virtud de (18.2.5), mostrar que para todo ideal maximal $\mathfrak m$ de $A$, $B_{\mathfrak m}/\mathfrak mB_{\mathfrak m}$ es una $(A/\mathfrak m)$-álgebra separable. Ahora bien, como $\mathfrak J$ está contenido en el radical de $A$ ($0_{\mathrm I}$, 7.1.10),

\oldpage[IV-4]{117}

se tiene $\mathfrak J\subset\mathfrak m$ y, si $\mathfrak m_0=\mathfrak m/\mathfrak J$, se tiene $A_0/\mathfrak m_0=A/\mathfrak m$ y $B_{\mathfrak m}/\mathfrak mB_{\mathfrak m}=(B_0)_{\mathfrak m_0}/\mathfrak m_0(B_0)_{\mathfrak m_0}$; la conclusión se sigue, por tanto, del hecho de que $B_0$ es una $A_0$-álgebra étale (18.2.5).

\begin{env}[18.3.3]
\label{IV.18.3.3}
\emph{Ejemplo.} --- La prop. (18.3.2) se aplica, en particular, cuando $A$ es un \emph{anillo local} separado y completo, siendo $\mathfrak J$ el ideal maximal de $A$, de modo que $A_0$ es un \emph{cuerpo} y la categoría de las $A_0$-álgebras finitas y étales es idéntica a la de las $A_0$-álgebras finitas y separables, que son, por tanto, isomorfas a productos directos de cuerpos, extensiones separables y finitas de $A_0$. En particular, si el cuerpo $A_0$ es \emph{separablemente cerrado}, estas extensiones son todas idénticas a $A_0$ y, por consiguiente, todo recubrimiento étale de $\Spec(A)$ es \emph{trivial} (18.2.7) en virtud de (18.3.2).
\end{env}

\begin{theorem}[18.3.4]
\label{IV.18.3.4}
Sean $A$ un anillo noetheriano y $\mathfrak J$ un ideal de $A$ tal que $A$ sea separado y completo para la topología $\mathfrak J$-preádica, $A_0=A/\mathfrak J$. Pongamos $S=\Spec(A)$, $S_0=\Spec(A_0)$. Sea $X$ un $S$-esquema propio sobre $S$, y pongamos $X_0=X\times_S S_0$. Entonces el funtor
\[
  Z\longmapsto Z\times_X X_0
\]
de la categoría de los $X$-esquemas finitos y étales sobre $X$ a la categoría de los $X_0$-esquemas finitos y étales sobre $X_0$ es una equivalencia de categorías.
\end{theorem}

Demostremos primero que este funtor es \emph{plenamente fiel}. Sean, pues, $Z'$ y $Z''$ dos $X$-esquemas finitos y étales sobre $X$. Pongamos $S_n=\Spec(A/\mathfrak J^{n+1})$, $X_n=X\times_S S_n$, $Z'_n=Z'\times_S S_n$, $Z''_n=Z''\times_S S_n$ para todo $n\geq 0$. Se sigue de (III, 5.4.1) que se tiene una biyección canónica $\Hom_X(Z',Z'')\xrightarrow{\sim}\varprojlim\Hom_{X_n}(Z'_n,Z''_n)$. Ahora bien, en virtud de (18.1.2), la aplicación canónica $\Hom_{X_{n+1}}(Z'_{n+1},Z''_{n+1})\to\Hom_{X_n}(Z'_n,Z''_n)$ es biyectiva, lo que termina de demostrar nuestra afirmación.

Queda por demostrar que, si $\sh{B}_0$ es una $\sh{O}_{X_0}$-álgebra finita y étale, existe una $\sh{O}_X$-álgebra finita y étale $\sh{B}$ y un isomorfismo $\sh{B}_0\xrightarrow{\sim}\sh{B}\otimes_{\sh{O}_X}\sh{O}_{X_0}$. Se sigue de (18.1.2) que existe un sistema proyectivo $(\sh{B}_n)$, donde $\sh{B}_n$ es una $\sh{O}_{X_n}$-álgebra finita y étale, y el segundo teorema de comparación (III, 5.1.4) demuestra la existencia de un $\sh{O}_X$-módulo coherente $\sh{B}$ y de un sistema proyectivo de isomorfismos $\sh{B}_n\xrightarrow{\sim}\sh{B}\otimes_{\sh{O}_X}\sh{O}_{X_n}$. El dato de una estructura de álgebra sobre un módulo $\sh{F}$ equivale al de un homomorfismo $\sh{O}_X\to\sh{F}$ que define la unidad y de un homomorfismo de multiplicación $\sh{F}\otimes_{\sh{O}_X}\sh{F}\to\sh{F}$ que hacen conmutar los diagramas de unidad, asociatividad y conmutatividad en los que solo intervienen potencias tensoriales de $\sh{F}$; se sigue de (III, 5.1.3), (I, 10.11.4) y (I, 10.11.7) que $\sh{B}$ está dotada de manera natural de una estructura de $\sh{O}_X$-álgebra para la cual el isomorfismo $\sh{B}_0\xrightarrow{\sim}\sh{B}\otimes_{\sh{O}_X}\sh{O}_{X_0}$ es un isomorfismo de álgebras. Además, $\sh{B}$ es un $\sh{O}_X$-módulo \emph{localmente libre}; esto se sigue también de (III, 5.1.3), (I, 10.11.4), (I, 10.11.7) y del hecho de que, en la categoría de los $\sh{O}_X$-módulos coherentes, los $\sh{O}_X$-módulos localmente libres $\sh{F}$ pueden definirse como aquellos para los que el funtor $\sh{G}\mapsto\mathscr{H}\!om_{\sh{O}_X}(\sh{F},\sh{G})$ es \emph{exacto}. Finalmente, para ver que $\sh{B}$ es una $\sh{O}_X$-álgebra étale, basta (18.2.5) mostrar que, para todo punto \emph{cerrado} $x\in X$, $\sh{B}_x\otimes_{\sh{O}_{X,x}}\kappa(x)$ es una $\kappa(x)$-álgebra separable. Pero como el morfismo estructural $f:X\to S$ es \emph{propio}, $f(x)$ es un punto cerrado de $S$ y, por tanto, pertenece a $S_0$, ya que $\mathfrak J$ está contenido en el radical de $A$ ($0_{\mathrm I}$, 7.1.10); la conclusión se sigue, por tanto, de que $X_0=f^{-1}(S_0)$ y de que $\sh{B}_0$ es una $\sh{O}_{X_0}$-álgebra finita y étale.
